% !TEX program = xelatex
\documentclass[12pt,a4paper]{article}

\usepackage{fontspec}
\usepackage{polyglossia}
\setmainlanguage{russian}
\setotherlanguage{english}
\setmainfont{DejaVu Serif}
\setsansfont{DejaVu Sans}
\setmonofont{DejaVu Sans Mono}

\usepackage{amsmath,amssymb,amsfonts,mathtools}
\usepackage{bm}
\usepackage{geometry}
\usepackage{enumitem}
\usepackage{hyperref}
\usepackage{microtype}
\usepackage{xcolor}
\geometry{margin=2.45cm}
\hypersetup{colorlinks=true,linkcolor=black,urlcolor=blue,citecolor=black}
\setlength{\emergencystretch}{3em}
\sloppy

\DeclareMathOperator{\Adm}{Adm}
\DeclareMathOperator{\Struct}{Struct}
\DeclareMathOperator{\Dist}{Dist}
\DeclareMathOperator{\Tr}{Tr}
\DeclareMathOperator{\Dom}{Dom}
\DeclareMathOperator{\Proj}{Proj}
\DeclareMathOperator{\Curv}{Curv}

\newcommand{\TMI}{\mathcal T_{\mathrm{MI}}}
\newcommand{\TOE}{\mathcal T_{\mathrm{MI-ToE}}}
\newcommand{\M}{\mathcal M_U}
\newcommand{\MI}{\mathcal M_I^{\mathrm{ToE}}}
\newcommand{\EI}{\mathcal E_I}
\newcommand{\AI}{\mathbb A_I}
\newcommand{\FI}{\mathbb F_I}
\newcommand{\Hilb}{\mathcal H}
\newcommand{\Prob}{\mathbb P}
\newcommand{\Lag}{\mathcal L}
\newcommand{\Unc}{\mathcal U}
\newcommand{\Eff}{\mathrm{eff}}
\newcommand{\SM}{\mathrm{SM}}
\newcommand{\GR}{\mathrm{GR}}
\newcommand{\QG}{\mathrm{QG}}
\newcommand{\dd}{\,d}

\newcommand{\statementbox}[1]{%
\begin{center}%
\fbox{\begin{minipage}{0.88\linewidth}\centering #1\end{minipage}}%
\end{center}%
}

\title{\textbf{TMI-ToE}\\[0.4em]
\Large Интерфейсно-полевая программа объединения фундаментальных взаимодействий\\
\large Новый проект в рамках Теории многообразных интерфейсов}
\author{Aleksey Salkutsan\\Independent Researcher}
\date{Версия: TMI-ToE-0.1 draft}

\begin{document}
\maketitle

\begin{abstract}
Данный документ фиксирует новый проект в рамках Теории многообразных интерфейсов: \emph{интерфейсно-полевую программу объединения фундаментальных взаимодействий}. Текст не утверждает, что готовая теория всего уже получена. Его статус осторожнее: после введения квантово-гравитационного слоя ТМИ допускает построение архитектуры, в которой гравитационное, электромагнитное, слабое и сильное взаимодействия рассматриваются как частные проекции единого физического интерфейсного многообразия.

Центральная идея проекта:
\statementbox{%
четыре фундаментальных взаимодействия являются различными проекциями единой интерфейсной кривизны.
}

Каноническая математическая основа ТМИ указывается как связанная работа:
\begin{center}
\href{https://doi.org/10.5281/zenodo.19978895}{Mathematical Abstraction of the Theory of Manifold Interfaces, DOI: 10.5281/zenodo.19978895}.
\end{center}
\end{abstract}

\section{Статус проекта}

Новый проект получает рабочее обозначение:
\[
\boxed{
\TOE
}
\]
и читается как \emph{TMI-ToE}: интерфейсно-полевая программа к объединённой теории фундаментальных взаимодействий.

Важно зафиксировать ограничение статуса:
\statementbox{%
Это не готовая теория всего и не утверждение о завершённом решении объединения фундаментальных взаимодействий. Это исследовательская программа, задающая формальный язык, в котором такое объединение можно строить осторожно и проверяемо.
}

Иначе говоря:
\statementbox{%
не готовая Theory of Everything, а interface-field program toward unification.
}

\section{Исходная мотивация}

В физическом и квантово-гравитационном слоях ТМИ уже была зафиксирована цепочка:
\statementbox{%
геометрия задаёт область возможного; квантовая физика задаёт распределение возможного; интерфейс переводит возможность в структуру.
}

Формально это выражалось как:
\[
 g
 \Rightarrow
 \Adm_c(g)
 \Rightarrow
 \Hilb_I(g)
 \Rightarrow
 \Psi_I
 \Rightarrow
 \Dist_{\mathrm{poss}}(\Psi_I;g)
 \Rightarrow
 \Struct_I .
\]

После осторожного перехода к квантовой гравитации сама геометрия перестаёт быть фиксированным фоном и становится квантовой возможностью:
\[
 g_{\mu\nu}
 \longrightarrow
 \widehat g_{\mu\nu},
 \qquad
 g\in\mathfrak G,
\]
где \(\mathfrak G\) обозначает пространство допустимых геометрий.

Тогда квантово-гравитационная цепочка принимает вид:
\[
\boxed{
\mathfrak G
\Rightarrow
\Psi_G[g]
\Rightarrow
\Prob_G(g)
\Rightarrow
I_G
\Rightarrow
 g_{\Eff}.
}
\]

Следующий естественный шаг: если гравитация тоже может быть описана в интерфейсной форме, то все четыре взаимодействия можно рассматривать как физические интерфейсы переноса, связи и структурирования состояний.

\section{Четыре фундаментальных взаимодействия как интерфейсы}

Вводится множество фундаментальных физических интерфейсов:
\[
\boxed{
\mathcal F_{\mathrm{phys}}
=
\{I_G,I_{EM},I_W,I_S\}.
}
\]

Здесь:
\begin{itemize}[leftmargin=2.2em]
\item \(I_G\) --- гравитационный интерфейс;
\item \(I_{EM}\) --- электромагнитный интерфейс;
\item \(I_W\) --- слабый интерфейс;
\item \(I_S\) --- сильный интерфейс.
\end{itemize}

Интерфейсная интерпретация взаимодействия:
\statementbox{%
фундаментальное взаимодействие есть интерфейсный канал преобразования физических состояний.
}

То есть взаимодействие понимается не только как сила или поле, а как допустимый физический канал, через который состояния одной структуры могут влиять на состояния другой структуры.

\section{Единое физическое интерфейсное многообразие}

Вводится центральный объект проекта:
\[
\boxed{
\MI
:=
(\M,\EI,\AI,\FI,\Psi,\chi_I,\Adm_I,\Struct_I).
}
\]

Где:
\begin{description}[leftmargin=3.2em,style=nextline]
\item[\(\M\)] пространство-время или базовое физическое многообразие;
\item[\(\EI\)] общее интерфейсное расслоение;
\item[\(\AI\)] универсальная интерфейсная связность;
\item[\(\FI\)] универсальная интерфейсная кривизна;
\item[\(\Psi\)] полное квантовое состояние физической системы;
\item[\(\chi_I\)] физическое интерфейсное поле;
\item[\(\Adm_I\)] множество допустимых физических переходов;
\item[\(\Struct_I\)] структурированное физическое состояние результата.
\end{description}

Главное утверждение проекта:
\[
\boxed{
\mathcal M_I^{\mathrm{ToE}}
\longrightarrow
\{I_G,I_{EM},I_W,I_S\}.
}
\]

То есть четыре взаимодействия рассматриваются как проекции единого физического интерфейсного многообразия:
\[
\boxed{
I_a=\pi_a(\MI),
\qquad
 a\in\{G,EM,W,S\}.
}
\]

\section{Универсальная интерфейсная связность}

Объединяющий объект физического уровня --- универсальная интерфейсная связность:
\[
\boxed{
\AI
=
\omega
+
A_{SU(3)}
+
A_{SU(2)}
+
A_{U(1)}
+
\chi_I.
}
\]

Здесь:
\begin{itemize}[leftmargin=2.2em]
\item \(\omega\) --- гравитационная или спиновая связность;
\item \(A_{SU(3)}\) --- калибровочное поле сильного взаимодействия;
\item \(A_{SU(2)}\) --- калибровочное поле слабого взаимодействия;
\item \(A_{U(1)}\) --- электромагнитно-гиперзарядный сектор;
\item \(\chi_I\) --- интерфейсное поле, связывающее геометрию, квантовое состояние и структурирование.
\end{itemize}

В более развёрнутой физической записи:
\[
\boxed{
\AI
=
\omega^{ab}J_{ab}
+
e^aP_a
+
G^A T_A
+
W^i T_i
+
B Y
+
\chi_I\Xi .
}
\]

Где:
\begin{itemize}[leftmargin=2.2em]
\item \(\omega^{ab}\) и \(e^a\) задают гравитационный сектор;
\item \(G^A\) --- глюонные поля сильного взаимодействия;
\item \(W^i\) --- слабые калибровочные поля;
\item \(B\) --- поле гиперзаряда;
\item \(J_{ab},P_a,T_A,T_i,Y\) --- соответствующие генераторы;
\item \(\Xi\) --- генератор или носитель интерфейсного сектора.
\end{itemize}

\section{Универсальная интерфейсная кривизна}

По аналогии с калибровочными теориями для \(\AI\) вводится кривизна:
\[
\boxed{
\FI
=
 d\AI
+
\AI\wedge\AI .
}
\]

Отдельные физические взаимодействия задаются проекциями этой кривизны:
\[
\boxed{
F_G=\pi_G(\FI),
\qquad
F_{EM}=\pi_{EM}(\FI),
}
\]
\[
\boxed{
F_W=\pi_W(\FI),
\qquad
F_S=\pi_S(\FI).
}
\]

Смысл этой записи:
\statementbox{%
Четыре фундаментальных взаимодействия являются четырьмя режимами или проекциями единой интерфейсной кривизны.
}

Это не доказывает объединение взаимодействий, но задаёт формальный язык, в котором такое объединение можно исследовать.

\section{Объединённое действие}

Минимальная архитектура действия проекта:
\[
\boxed{
S_{\mathrm{TMI-ToE}}
=
S_{\GR}
+
S_{\SM}
+
S_I.
}
\]

Развёрнуто:
\[
\boxed{
S_{\mathrm{TMI-ToE}}
=
\int_{\M}
\sqrt{-g}
\left[
\frac{c^4}{16\pi G}R
+
\Lag_{\SM}
+
\Lag_I(\chi_I,\nabla\chi_I,\AI,\FI,\Psi)
\right]
\dd^4x .
}
\]

Где:
\begin{itemize}[leftmargin=2.2em]
\item \(S_{\GR}\) --- гравитационная часть;
\item \(S_{\SM}\) --- действие Стандартной модели;
\item \(S_I\) --- интерфейсный член действия;
\item \(\Lag_I\) --- лагранжиан интерфейсного поля.
\end{itemize}

Именно \(\Lag_I\) должен отвечать за:
\begin{enumerate}[leftmargin=2.2em]
\item переход от распределения возможностей к структуре;
\item обратную связь между геометрией и квантовым состоянием;
\item связь гравитационного и калибровочных секторов;
\item возможные интерфейсные поправки к декогеренции, геометрии и вероятностям переходов.
\end{enumerate}

\section{Уравнения поля}

Из действия следуют три класса вариационных уравнений:
\[
\boxed{
\frac{\delta S_{\mathrm{TMI-ToE}}}{\delta g^{\mu\nu}}=0,
\qquad
\frac{\delta S_{\mathrm{TMI-ToE}}}{\delta \Psi}=0,
\qquad
\frac{\delta S_{\mathrm{TMI-ToE}}}{\delta \chi_I}=0.
}
\]

Первое уравнение задаёт гравитационную динамику, второе --- квантовую динамику, третье --- динамику интерфейсного поля.

В предельной эффективной форме гравитационный сектор может быть записан как:
\[
\boxed{
G_{\mu\nu}[g]
=
\frac{8\pi G}{c^4}
\left(
T_{\mu\nu}^{(\Psi)}
+
T_{\mu\nu}^{(I)}
\right),
}
\]
где \(T_{\mu\nu}^{(I)}\) --- вклад интерфейсного слоя.

\section{Квантовый слой: от возможности к структуре}

Квантовое состояние задаёт распределение возможностей, а интерфейсное поле задаёт структурную актуализацию.

До интерфейсного перехода состояние может быть записано как:
\[
|\Psi\rangle
=
\sum_\alpha c_\alpha |s_\alpha\rangle .
\]

В плотностной записи:
\[
\rho
=
|\Psi\rangle\langle\Psi|
=
\sum_{\alpha,\beta}
 c_\alpha\overline{c_\beta}
 |s_\alpha\rangle\langle s_\beta| .
\]

Интерфейсный переход можно представить как квантовый канал:
\[
\boxed{
\mathcal E_I(\rho)
=
\sum_k K_k^{(I)}\rho K_k^{(I)\dagger},
\qquad
\sum_k K_k^{(I)\dagger}K_k^{(I)}=\mathbf 1.
}
\]

После интерфейсного структурирования:
\[
\boxed{
\rho
\overset{I}{\longmapsto}
\rho_I'
=
\mathcal E_I(\rho)
\overset{\mathfrak S_I}{\longmapsto}
\Struct_I.
}
\]

В пределе полной редукции:
\[
\rho_I'
\approx
\sum_k p_k |s_k\rangle\langle s_k|,
\qquad
p_k=\Tr(E_k^{(I)}\rho),
\qquad
E_k^{(I)}=K_k^{(I)\dagger}K_k^{(I)}.
\]

Осторожная физическая формулировка:
\statementbox{%
Интерфейс не объявляется метафизическим разрушителем суперпозиции. Он задаёт механизм структуризации, который может быть выражен как декогеренция, редукция, регистрация результата или квантовый канал относительно выбранной структуры различения.
}

\section{Главная цепочка TMI-ToE}

Объединённая цепочка проекта:
\[
\boxed{
\AI
\Rightarrow
\FI
\Rightarrow
\{F_G,F_{EM},F_W,F_S\}
\Rightarrow
\Adm_I
\Rightarrow
\Hilb_I
\Rightarrow
\Psi
\Rightarrow
\Dist_{\mathrm{poss}}
\Rightarrow
\Struct_I .
}
\]

Смысл:
\begin{enumerate}[leftmargin=2.2em]
\item универсальная интерфейсная связность задаёт объединённый физический канал;
\item универсальная интерфейсная кривизна задаёт режимы взаимодействия;
\item четыре фундаментальных взаимодействия являются проекциями этой кривизны;
\item допустимые переходы задают область физически возможного;
\item квантовое состояние задаёт распределение возможностей;
\item интерфейсное поле переводит распределение возможностей в структуру.
\end{enumerate}

В максимально сжатой форме:
\statementbox{%
фундаментальная физика есть динамика единого интерфейсного многообразия.
}

\section{Предельные случаи}

Чтобы программа могла претендовать на физическую содержательность, она должна восстанавливать известные теории как предельные случаи.

\subsection{Предел общей теории относительности}

Если интерфейсный вклад исчезает или становится пренебрежимо малым:
\[
T_{\mu\nu}^{(I)}\to 0,
\]
то должно получаться стандартное уравнение Эйнштейна:
\[
\boxed{
G_{\mu\nu}[g]
=
\frac{8\pi G}{c^4}T_{\mu\nu}.
}
\]

\subsection{Предел квантовой механики}

Если геометрия фиксирована:
\[
 g_{\mu\nu}=g_{\mathrm{fixed}},
\]
то должна восстанавливаться обычная квантовая динамика:
\[
\boxed{
 i\hbar\partial_t\Psi=\widehat H\Psi.
}
\]

\subsection{Предел Стандартной модели}

Если гравитационный и интерфейсный секторы отделяются, должна восстанавливаться калибровочная группа Стандартной модели:
\[
\boxed{
SU(3)\times SU(2)\times U(1).
}
\]

При этом должны сохраняться известные поля, заряды, массы и калибровочные взаимодействия.

\section{Что ещё не доказано}

Проект TMI-ToE на данном этапе не доказывает:
\begin{itemize}[leftmargin=2.2em]
\item происхождение всех частиц Стандартной модели;
\item значения масс и зарядов;
\item число поколений фермионов;
\item строгую квантовую меру на пространстве геометрий;
\item уникальный лагранжиан \(\Lag_I\);
\item проверенное экспериментальное предсказание.
\end{itemize}

Поэтому корректный статус:
\statementbox{%
архитектура кандидата на объединённую теорию, а не завершённая теория всего.
}

\section{Минимальные условия перехода к физической теории}

Для перехода от программы к физической теории нужно выполнить следующие задачи:
\begin{enumerate}[leftmargin=2.2em]
\item построить конкретный лагранжиан \(\Lag_I\);
\item получить из него уравнение динамики \(\chi_I\);
\item показать восстановление ОТО, КМ и Стандартной модели;
\item вывести хотя бы одну вычислимую интерфейсную поправку;
\item предложить наблюдаемую величину или экспериментальный критерий.
\end{enumerate}

Минимальная вычислимая величина-кандидат:
\[
\boxed{
|\rho_{ij}(t)|
=
\exp\left(-\int_0^t\Gamma_I(s)\dd s\right)
|\rho_{ij}(0)|,
}
\]
где \(\Gamma_I\) --- интерфейсная скорость декогеренции.

Если \(\Gamma_I\) удаётся выразить через \(\chi_I\), кривизну \(\FI\) и квантовое состояние \(\Psi\), то проект получает первую вычислимую физическую модель.

\section{Дорожная карта проекта}

\subsection*{TMI-ToE-0.1}
Фиксация архитектуры:
\[
\AI,
\quad
\FI,
\quad
\pi_a(\FI),
\quad
S_{\mathrm{TMI-ToE}},
\quad
\Lag_I.
\]

\subsection*{TMI-ToE-0.2}
Уточнение математической природы интерфейсного расслоения \(\EI\), универсальной связности \(\AI\) и проекций \(\pi_a\).

\subsection*{TMI-ToE-0.3}
Построение первого toy model для \(\Lag_I\) и \(\Gamma_I\).

\subsection*{TMI-ToE-0.4}
Проверка предельных случаев:
\[
\TOE\to \GR,
\qquad
\TOE\to \mathrm{QM},
\qquad
\TOE\to \SM.
\]

\subsection*{TMI-ToE-1.0}
Публикационная версия как \emph{Interface-Field Program Toward a Unified Theory of Fundamental Interactions}.

\section{Итоговая формулировка}

Итоговый тезис нового проекта:
\statementbox{%
После введения квантово-гравитационного слоя Теория многообразных интерфейсов допускает расширение до интерфейсно-полевой программы объединения фундаментальных взаимодействий. В этой программе гравитационное, электромагнитное, слабое и сильное взаимодействия рассматриваются как различные проекции единой интерфейсной кривизны, а физическая реальность описывается как динамика единого интерфейсного многообразия.
}

Главная формула:
\[
\boxed{
\pi_G(\FI)=F_G,
\qquad
\pi_{EM}(\FI)=F_{EM},
\qquad
\pi_W(\FI)=F_W,
\qquad
\pi_S(\FI)=F_S.
}
\]

Сверхсжатая форма:
\statementbox{%
четыре взаимодействия = четыре проекции единой интерфейсной кривизны.
}

\end{document}
